\scanpage{66}
\mspage{II. 14}
% lines 1-12
En effet, le groupe algébrique \inserted{$G\subset\operatorname{Aut}_{\mathbb{Q}_\ell}(\unsure{E_\ell|E})$} engendré par cette représentation est tel que $G^\circ$ est nilpotent, \deleted{Donc} \deleted{\unsure{par}} de \inserted{Donc :} un \uline{produit} \inserted{\uline{direct}} $T\times U$, $T$ un tore, $U$ unipotent. \inserted{\uline{A prouver} $U=1$.} Soit $\pi'\subset\pi$ l'image inverse de $G^\circ$; il correspond à un rev. \deleted{fini} étale connexe $S'$ de $S$, et quitte à passer à celui-ci, on peut supposer $G=G^\circ$. On utilise le th. de \uline{\unsure{Čebotareff|Čebotaref}} que \uline{les éléments de frobénius} \uline{sont denses} \uline{dans} $\pi_1(S,\unsure{\xi|\zeta})$. On en conclut que leurs images $f_s$ dans $G$ \uline{engendrent} $G$ comme groupe algébrique. Or $T$ est \unsure{l'ens} des élts semi-simples de $G$, donc $G=T$, i.e. $U=(1)$, cqfd.

% lines 13-19
b) \uline{Fin de la démonstration}. Il reste à prouver que \deleted{\unsure{pour tt}} pour tt $\mathfrak m\in\operatorname{Max}(R)$ tel que $R/\mathfrak m$ de car. $\ell$, $\exists$ extension finie $K'$ de $K$ telle que $A_{\deleted{\unsure{K}}K'}$ ait bonne réduction en les idéaux maximaux de $V'$ (où $V'$ est le normalisé de $V=R_{\mathfrak m}$ dans $K'$). Pour ceci, \inserted{on note que \deleted{l'action} en vertu \unsure{de|du} cor. 2, l'action de $\unsure{\operatorname{Gal}|\pi\mathrel{\text{.}}\operatorname{Gal}}(\overline K/K)$ sur $T_\ell(A(\overline K))\otimes_{\mathbb{Z}_\ell}\mathbb{Q}_\ell$ est $\frac12$ simple abélienne, donc il en est de \unsure{même}; \unsure{a fortiori} de celle du gr\textsuperscript{t} d'inertie $I_{\mathfrak m}$.}

% cancelled paragraph, lower block crossed diagonally
\deleted{\deleted{utilisant} un résultat d'un exposé ultérieur, on utilise le fait que \deleted{\unsure{il existe}} l'on peut trouver une extension \deleted{de} finie $K'/K$ tel que \deleted{le schéma} \inserted{le modèle de Néron de} $A_{K'}$ sur $R'$ ait ``réduction $\frac12$ st.R'' en tout point de \deleted{\unsure{$V'$}} $\operatorname{Spec}V'$. \deleted{Alors} \unsure{bien plus,}}
% archive sheet number: 65
